\documentclass[fontset=none,12pt,a4paper]{ctexart}
\usepackage[margin=2.4cm]{geometry}
\usepackage{fontspec}
\usepackage{xcolor}
\usepackage{enumitem}
\usepackage{tcolorbox}
\usepackage{titlesec}
\usepackage{parskip}

\IfFontExistsTF{Songti SC}
  {\setCJKmainfont{Songti SC}\setCJKsansfont{Songti SC}\setCJKmonofont{Songti SC}}
  {\setCJKmainfont[BoldFont=FandolSong-Bold.otf]{FandolSong-Regular.otf}\setCJKsansfont[BoldFont=FandolSong-Bold.otf]{FandolSong-Regular.otf}\setCJKmonofont[BoldFont=FandolSong-Bold.otf]{FandolSong-Regular.otf}}
\IfFontExistsTF{TeX Gyre Termes}{\setmainfont{TeX Gyre Termes}}{\setmainfont{Times New Roman}}

\definecolor{RuleBlue}{HTML}{EAF3FF}
\definecolor{RuleLine}{HTML}{5C88C5}
\definecolor{RuleGreen}{HTML}{EEF8EF}
\definecolor{RuleGreenLine}{HTML}{5B9B6A}

\titleformat{\section}{\Large\bfseries}{\thesection.}{0.6em}{}
\setlength{\parindent}{0pt}
\setlist[enumerate]{leftmargin=2.2em,itemsep=0.7em}
\setlist[itemize]{leftmargin=2.0em,itemsep=0.35em}

\newcommand{\markexample}[1]{\texttt{#1}}

\title{\textbf{参与者作答规则}}
\author{}
\date{}

\begin{document}
\maketitle
\vspace{-2em}

\begin{tcolorbox}[colback=RuleBlue,colframe=RuleLine,title=实验安排]
请按照任务清单依次完成材料。根据实验计划，每位参与者会完成：
\begin{itemize}
  \item 1 个 \textbf{Human-only} packet：只依据题面材料自行判断。
  \item 1 个 \textbf{AI辅助} packet：题面中已经给出 AI 初稿，可接受或修改。
\end{itemize}
两种条件使用同样的 case packet 类型和同样的最终评分目标：给每一句话标出句子级月份范围。
\end{tcolorbox}

\section*{Human-only 作答规则}
\begin{tcolorbox}[colback=white,colframe=black!25,title=Human-only]
Human-only 条件下，PDF 只提供目标文本和答题表，不提供 AI 初稿。你的任务是：
\begin{enumerate}
  \item 给目标段中的每一句标出起始时间和结束时间。
  \item 如果一句话只对应一个时间点或一个月份，起始时间和结束时间可以写成同一个时间。
  \item 如果无法确定时间范围，请标为“未知”，并可在备注处简单说明原因。
\end{enumerate}
不要填写 TimeBlock 顺序、时间标志物补全、ISO 日期等中间变量；这里只看最终句子级时间范围。
\end{tcolorbox}

\section*{AI辅助作答规则}
\begin{tcolorbox}[colback=gray!6,colframe=black!20,title=AI辅助]
AI辅助条件下，PDF 中已经预先给出 AI 草稿。你可以：
\begin{enumerate}
  \item 直接接受 AI 的时间范围。
  \item 修改开始年月或结束年月。
  \item 把该句改成“未知”。
  \item 按你的判断修正 AI 给出的下沉、插叙或备注信息。
\end{enumerate}

AI 预标的下沉句和插叙句会在句子正文前显示
\markexample{【下沉】} 或 \markexample{【插叙】}。AI 给出的跨文本 TimeBlock 来源、下沉原因等备注，会在对应句子下方以单独的“备注”行显示。

其中“跨文本 TimeBlock”来自 AI 对另一篇纪传中相近时间块的自动对齐，例如
\markexample{跨文本TimeBlock：8.23.5} 表示当前句子所在时间块参考了以 8.23.5 开头的另一文本时间块；它不是句子一一对应标记，只作为时间判断参考。
\end{tcolorbox}

\section*{共同时间标注规则}
\begin{tcolorbox}[colback=white,colframe=black!25,title=时间范围]
两种条件都请给每一句话划分时间范围，并在句子序号前标记：

\vspace{0.4em}
\centerline{\Large\markexample{（起始时间，结束时间）}}
\vspace{0.4em}

如果一句话只对应一个时间点或一个月份，起始时间和结束时间可以写成同一个时间。
\end{tcolorbox}

\begin{tcolorbox}[colback=RuleGreen,colframe=RuleGreenLine,title=纪年规则]
\textbf{注意：十月为一年的开始。}

遇到汉初相关纪年时，请按题面材料和上下文判断年份；不要默认正月是一年的开始。
\end{tcolorbox}

\section*{特殊情况}
\begin{enumerate}
  \item \textbf{插叙}

  如果句子是插叙、追述或倒叙，请在句子序号前标记“插叙”，并继续填写时间范围：

  \vspace{0.2em}
  \centerline{\Large\markexample{插叙（起始时间，结束时间）}}

  \item \textbf{下沉}

  如果句子没有具体事件发生，只是人物身份、背景、评价、说明来源等描述性语句，请标记为：

  \vspace{0.2em}
  \centerline{\Large\markexample{下沉}}
\end{enumerate}

\vfill
\begin{tcolorbox}[colback=gray!6,colframe=black!20]
作答时请优先依据当前 PDF 中给出的材料与上下文判断；无法确定时，可在备注处说明不确定原因。
\end{tcolorbox}

\end{document}
